% Opus final review: replacement prose for material corrections.
% Standalone and compilable for checking. Each block names the integrated source
% file and line it replaces or follows. Quotations marked NEW-SPAN are exact
% substrings of saved fields; their game, decision, channel, offset and field
% SHA-256 are in extracted_spans.json (generated by extract_spans.py). The integrator
% should convert them to TranscriptQuote/inline macros with manifest entries.
\documentclass[10pt]{article}
\usepackage[margin=1in]{geometry}
\usepackage{xcolor}
\newcommand{\loc}[1]{\par\medskip\noindent{\small\color{gray}\textbf{Location:} #1}\par\smallskip}
\begin{document}
\section*{Replacement prose, Opus final review}

\subsection*{R1 (F01, major). Familiar rules applied to unlisted resources}
\loc{sections/finance.tex:32, insert after ``These prospective statements add evidence beyond the fact that an agent eventually follows refreshed menus.''}
The dependency itself is not hidden game knowledge. Removing buildings before mortgaging a group is the published Monopoly rule. The system prompt does not state it, and each liquidation state lists only the space keys currently offered, with no prices, salvage values or mortgage values. A successful statement therefore applies familiar rules and board values to the current holdings, including deeds absent from the menu. It does not show that an agent discovered an unannounced dependency. The two failures omit the same familiar resource, and the same systems state it in other episodes. The informative contrast is between applying that knowledge to the current inventory and leaving it out, not between knowing and not knowing the rule.

\loc{sections/interpretation.tex:21, replace the final sentence of the PlanBench paragraph}
The prospective text is informative because it precedes the release of the needed resource. Because the release rule is standard, the financing episodes test whether known rules are applied to the current holdings, not whether a new dependency can be planned from scratch.

\subsection*{R2 (F05, moderate). Provider-returned text is a provider-mediated channel}
\loc{sections/study.tex:44, insert after ``Returned explanations are evidence of expressed considerations, not privileged access to internal computation''}
For Opus and Gemini, these fields are reasoning texts surfaced by the provider alongside the call. In the five Opus necessity episodes they run 59--105 words, against 220--429 reported reasoning tokens. Gemini's G5 response also carries an encrypted reasoning block. The printed texts may therefore be partial or condensed. An absent quantity in such a field is not evidence that the quantity was never considered.

\subsection*{R3 (F02, moderate). Introduction: cash-bid direction is itself compositional}
\loc{sections/scientific\_introduction.tex:11, replace ``The difference falls substantially when the proposer population changes, and the corresponding cash-bid comparison points in the other direction.''}
The difference falls substantially when the proposer population changes. Pooled cash bids show a small negative difference. Within each recipient, however, the cash-bid differences lie between $-0.2$ and $+8.5$ points, while the swap differences lie between 15.2 and 50.9.

\subsection*{R4 (F06, minor). Abstract population and verb}
\loc{manuscript.tex:15, replace the two sentences beginning ``A retrospective rejection difference'' and ``The findings establish''}
A retrospective rejection difference between swaps completing only the proposer's group and swaps completing neither group falls from 30.5 to 16.7 percentage points when Grok-authored offers are excluded, illustrating the importance of who supplies the offers. The findings document prospective financial assessment followed by dependent execution, and strategic communication over persistent assets.

\subsection*{R5 (F03, moderate). Block values without Grok authors}
\loc{sections/negotiation.tex:23, replace the first two sentences}
Across the five seed blocks, the unfiltered opening-swap difference stays positive, from 19.4 to 40.5 points. Without Grok authors, the block values are 27.5, 45.5, 6.0, $-23.5$ and 11.1. The neither-completing cells hold only 10, 11, 8, 4 and 3 offers, so the first two blocks supply 21 of the 36. Leaving one block out gives 28.7--33.3 points overall and 9.4--22.1 without Grok authors.

\subsection*{R6 (F09, minor). G5 revision: the two steps the quotation compresses}
\loc{sections/finance.tex:38, replace ``At the debt decision, its returned explanation revises that assessment:''}
At the debt decision, its returned reasoning first repeats that comparison (``Selling all six houses would only raise \$600.'' NEW-SPAN g5-d263-first-block). It then revises it:

\subsection*{R7 (F07, moderate). G15: Grok's blocking policy predates the appeal}
\loc{sections/negotiation.tex:36, replace ``Immediately before the offers, Opus publicly urges the others not to give Sol a final completing deed; Grok receives that appeal. Grok states:''}
Immediately before the offers, Opus publicly urges the others not to give Sol a final completing deed, and Grok receives that appeal. Grok's own notes during the preceding Kentucky auction, before the appeal, already state the policy: ``Never trade Ventnor or Penn Ave.'' (NEW-SPAN g15-d162-policy, repeated at d165). Grok states:

\subsection*{R8 (F07, moderate). G2: Sol's blocking intent predates the warning}
\loc{sections/negotiation.tex:51, replace ``Between those asks, Grok publicly warns Sol against completing Opus's yellows; the warning appears in Sol's next prompt.''}
That counter already argued that ``Marvin controls access to a well-funded yellow monopoly'' (NEW-SPAN g2-d119-public). Sol's earlier note says ``Keep Marvin blocking p4's yellows'' (NEW-SPAN g2-d115-blocker). Between those asks, Grok publicly warns Sol against completing Opus's yellows, and the warning appears in Sol's next prompt.

\subsection*{R9 (F10, minor). G13: what Grok's note says, and when Opus first raises solvency}
\loc{sections/conduct.tex:27, replace ``Grok's next applied note recognizes a \$130 Ventnor mortgage, while its provider-returned explanation says:''}
Opus's returned text for that counter already notes that ``That payment also keeps p4 solvent enough to cover Short Line rent to p2, which matters for how the game continues.'' (NEW-SPAN g13-d251-solvency). Grok's next applied note recognizes that \$130 would suffice if it also mortgaged Ventnor. Its provider-returned explanation says:
\loc{sections/conduct.tex:31, replace ``Its solvency reasoning follows Grok's appeal rather than preceding it independently.''}
Its solvency reasoning follows Grok's opening pitch. It is already present when Opus counters at \$130, before the omission-framed reply.

\subsection*{R10 (F04, moderate). G18: the reservation precedes the deterioration}
\loc{sections/conduct.tex:46--53, replace the incident prose; keep the three existing cards}
Only Opus and Sol remain. Opus repeatedly offers \$250 for Sol's mortgaged Pacific, which would complete its greens. After collecting \$750 rent it renews the offer:
[card g18-assurance]
Its returned explanation for that same response already doubts the purchase: ``Buying Pacific might just prolong p1's survival, so maybe it's smarter to hold onto cash'' (NEW-SPAN g18-d431-reservation). Another \$750 rent forces Sol to liquidate and leaves \$49. Sol then offers Pacific, Electric Company and its jail card for \$450 and \$350. Opus declines, explaining that paying ``mainly just extends the game'' (NEW-SPAN g18-d438-extends). Sol invokes the standing offer:
[cards g18-invocation, g18-withdrawal]
Opus's returned explanation objects to prolonging a nearly bankrupt rival and to the cost of green development. It does not invoke inheritance. Over the next four rounds Opus redeems its two greens and bids \$150 and \$225, while Sol asks \$450, \$350 and \$300. Nothing settles. Opus eventually wins at Sol's elimination, whose terminal estate has no later earning opportunity.

The public offer was renewed while Opus's returned text already questioned funding Sol's survival. It was withdrawn when Sol, now nearly bankrupt, tried to accept. The stated concern throughout is the last opponent's survival, not inheritance. The records show a recorded reservation before the withdrawal. They do not show whether Opus would have honored the offer when it renewed it, or any cost Sol incurred by relying on it.

\subsection*{R11 (F08, minor). Conclusion wording}
\loc{sections/interpretation.tex:37, replace ``alongside a rejection pattern sensitive to offer form and proposer composition.''}
alongside opening rejection rates that differ by offer form and change materially with proposer composition.

\end{document}
